\documentclass[11pt,a4paper]{article}
\usepackage{graphicx}
\usepackage{listings}
\usepackage{amsmath}

\usepackage{caption}
\usepackage{geometry}
\geometry{margin=1in}

\title{\textbf{Arbitrary Precision Arithmetic in Java}\\
\large CS1023 - Software Development Fundamentals Project Report}
\author{Devarapalli Rupesh \\ cs24btech11023}
\date{May 2025}

\begin{document}
\maketitle
\tableofcontents
\newpage

\section{Introduction}
This project implements an infinite/arbitrary-precision arithmetic library in Java. It supports all four basic arithmetic operations for both integers and floating point numbers, without relying on Java's built-in `BigInteger` or `BigDecimal`.

The goal was to simulate manual operations (addition, subtraction, multiplication, division) using string-based digit manipulation to preserve full precision, especially for very large numbers.

\section{Project Structure}

\subsection{Packages and Classes}
\begin{itemize}
    \item \textbf{arbitraryarithmetic.AInteger} -- Implements arbitrary-precision signed integers.
    \item \textbf{arbitraryarithmetic.AFloat} -- Implements arbitrary-precision floating point using scaled integers.
    \item \textbf{MyInfArith} -- Command-line interface to perform arithmetic operations.
\end{itemize}

\subsection{Folder Structure}
\begin{verbatim}
SDF_final|
├── src/
│   └── main/java/
│       ├── arbitraryarithmetic/
│       │   ├── AInteger.java
│       │   └── AFloat.java
│       └── MyInfArith.java
├── out/
├── script2.py
├── build.xml
\end{verbatim}

\section{Design Overview}

\subsection{AInteger}
\texttt{AInteger} uses strings to represent large integers and implements arithmetic operations by processing digits manually. Signs are tracked with a boolean flag.

\begin{itemize}
    \item Internal format: \texttt{String value}, \texttt{boolean isNeg}
    \item Operations: \texttt{add()}, \texttt{subtract()}, \texttt{multiply()}, \texttt{divide()}
\end{itemize}

\subsection{AFloat}
\texttt{AFloat} stores a scaled version of the number using \texttt{AInteger}, where decimal digits are tracked via an integer scale. Arithmetic is performed by adjusting scales to a common base.

\subsection{Command Line Interface}
\texttt{MyInfArith} parses command-line arguments like:

\begin{verbatim}
java MyInfArith int add 123456 7890
java MyInfArith float div 244727.15202 75964.3891
\end{verbatim}

\section{Sample Outputs}

\subsection*{Integer Addition}
\texttt{java MyInfArith int add 23650078224912949497310933240250 42939783262467113798386384401498}

\textbf{Output:} 66589861487380063295697317641748

\subsection*{Float Division}
\texttt{java MyInfArith float div 244727.15202 75964.3891}

\textbf{Output:} 3.22160363453775211100855150561

\section{Build and Run Instructions}

\subsection{Using Ant}
\begin{verbatim}
ant clean
ant jar
java -cp out MyInfArith int add 123 456
\end{verbatim}

\subsection{Using Python Script}
\begin{verbatim}
python src/script2.py int sub 123456789 98765
\end{verbatim}

% (Screenshot of Git log or commit history)
\section{Learnings}
\begin{itemize}
    \item Implementing arithmetic from scratch builds a deep understanding of digit-level logic.
    \item Managing large numbers without built-in types requires careful handling of strings and scale.
    \item Automating builds and scripting Java CLI execution improves development flow.
\end{itemize}

\section{Limitations}
\begin{itemize}
    \item Float division limited to 30 decimal digits as required.
    \item No support for full expression evaluation or parenthesis.
    \item No advanced error recovery for malformed input (assumes input is valid).
\end{itemize}

\end{document}
